\section{Pre-registration and analysis log}
\label{app:prereg}

Every confirmatory family was written into a protocol file before the models it concerns were trained or the numbers it
concerns were computed; the SHA-256 hash of the file and the time were recorded at that moment (\tableref{tab:app-protocols}).
Later edits are listed as amendments with their own hashes. \tableref{tab:app-tests} lists every pre-specified contrast
with its expectation and outcome. Four expectations were not met (TGR against the source positive rate for ResNet-18 at
the natural prevalence, TENT against KD + TGA, TGR against the source cutoff for ConvNeXt-T at $\pi=0.2$, and
non-inferiority to the seed-2 training of H-optimus-0, F2), two were met only in part (TENT against AdaBN) and one
result was better than expected (N2d).
The ResNet-18 arms with the H-optimus-0 teacher (no-teacher student, KD, KD + TGA, no-teacher adaptation) were run first
and analysed before these protocols; the expectations of the later families were derived from them.

Analyses not covered by a protocol are descriptive and labelled as such: the attribution and the share of the teacher's
advantage (\figureref{fig:gap}), the additional recalibration baselines (UPA, one bias to the source mean, EM with
bias-corrected calibration, BBSE, bias + temperature and quantile mapping to the teacher), the labelled few-image cutoff,
the prevalence estimates, teacher test-time augmentation and the throughput measurements. The additional baselines
were added after a literature review, once all pre-specified results of Protocols R and P had been computed; they do
not change any pre-specified result.

\begin{table}[htbp]
\floatconts{tab:app-protocols}
{\caption{Protocols in the order they were frozen (UTC). Each protocol file states the families, the expected direction
of every contrast and what was already known when it was written.}}
{%
\scriptsize
\setlength{\tabcolsep}{4pt}
\begin{tabular}{@{}llp{7.2cm}@{}}
\toprule
Protocol & Frozen & Content \\
\midrule
E1.7 (multi-teacher) & 2026-11-25 08:42 & Virchow2, UNI, UNI2-h teachers and ResNet-18 students; 9 contrasts. Amendment 1 (10:44): ConvNeXt-T students with H-optimus-0; amendment 2 (12:04): teacher LoRA seeds 1--2 (descriptive). \\
R (recalibration) & 2026-11-25 08:55 & TGR against label-free baselines (9 contrasts) and against AdaBN with 1, 2, 5 images (3 contrasts). \\
T (TENT) & 2026-11-25 10:40 & TENT against AdaBN and against retraining, KD and no-teacher students (4 contrasts). \\
P (prevalence shift) & 2026-11-25 18:39 & Resampling design, TGR against the source positive rate and the source cutoff at $\pi=0.2, 0.6$ (8 contrasts). \\
N (non-inferiority) & 2026-11-25 22:03 & Margin 0.02; ConvNeXt-T against H-optimus-0 (2 contrasts) and Virchow2 (superiority + 2 contrasts). \\
F and C & 2026-11-26 06:44 & Teacher-seed robustness (2 contrasts + key secondary) and the ConvNeXt-T no-teacher control (2 contrasts + wording rule). \\
\bottomrule
\end{tabular}
}
\end{table}

\begin{table}[htbp]
\floatconts{tab:app-tests}
{\caption{Every pre-specified contrast, in the order the protocols were written. R18 / CNX: ResNet-18 / ConvNeXt-T student. Differences are new-scanner F1 with 95\% paired-bootstrap intervals; $p$ is Holm-adjusted within the family (one-sided for non-inferiority, N1a--N2e and F; unadjusted for the key secondary estimate of F, which is outside the Holm family and includes the already-known seed-0 term). NI: non-inferiority at the margin 0.02.}}
{%
\tiny
\setlength{\tabcolsep}{2.5pt}
\begin{tabular}{@{}l>{\raggedright\arraybackslash}p{5.0cm}ccp{1.8cm}>{\raggedright\arraybackslash}p{2.4cm}@{}}
\toprule
Family & Contrast & Difference [95\% CI] & Holm $p$ & Expected & Outcome \\
\midrule
E1.7 & R18, Virchow2: KD $-$ no-teacher & +0.010 [-0.009, +0.030] & 1.000 & small, n.s. & met \\
E1.7 & R18, Virchow2: KD+TGA $-$ KD & +0.060 [+0.043, +0.079] & $<$0.001 & $>0$, significant & met \\
E1.7 & R18, Virchow2: KD+TGA $-$ no-teacher adaptation & +0.012 [-0.001, +0.027] & 0.367 & 0 to +0.03 & met \\
E1.7 & R18, UNI: KD $-$ no-teacher & +0.005 [-0.018, +0.027] & 1.000 & small, n.s. & met \\
E1.7 & R18, UNI: KD+TGA $-$ KD & +0.066 [+0.046, +0.091] & $<$0.001 & $>0$, significant & met \\
E1.7 & R18, UNI: KD+TGA $-$ no-teacher adaptation & +0.013 [+0.001, +0.027] & 0.227 & 0 to +0.03 & met \\
E1.7 & R18, UNI2-h: KD $-$ no-teacher & -0.007 [-0.031, +0.015] & 1.000 & small, n.s. & met \\
E1.7 & R18, UNI2-h: KD+TGA $-$ KD & +0.067 [+0.044, +0.094] & $<$0.001 & $>0$, significant & met \\
E1.7 & R18, UNI2-h: KD+TGA $-$ no-teacher adaptation & +0.002 [-0.014, +0.019] & 1.000 & 0 to +0.03 & met \\
R1 & R18 KD: TGR $-$ source cutoff & +0.037 [+0.022, +0.053] & $<$0.001 & $>0$ & met \\
R1 & R18 KD: TGR $-$ source positive rate & +0.002 [-0.008, +0.012] & 0.675 & $>0$ & not met \\
R1 & R18 KD: TGR $-$ EM prior & +0.210 [+0.176, +0.240] & $<$0.001 & $>0$ & met \\
R1 & R18 KD: TGR $-$ teacher pseudo-label cutoff & +0.004 [-0.001, +0.011] & 0.400 & none & -- \\
R1 & R18 KD: TGR $-$ AdaBN & -0.010 [-0.024, +0.004] & 0.400 & $<0$ & direction met (n.s.) \\
R1 & CNX KD: TGR $-$ source cutoff & +0.036 [+0.022, +0.052] & $<$0.001 & $>0$ & met \\
R1 & CNX KD: TGR $-$ source positive rate & +0.019 [+0.009, +0.031] & $<$0.001 & $>0$ & met \\
R1 & CNX KD: TGR $-$ EM prior & +0.129 [+0.098, +0.161] & $<$0.001 & $>0$ & met \\
R1 & CNX KD: TGR $-$ teacher pseudo-label cutoff & +0.006 [+0.002, +0.012] & 0.022 & none & -- \\
R2 & R18 KD, 1 image: TGR $-$ AdaBN & +0.025 [+0.010, +0.041] & 0.004 & $\ge 0$ & met \\
R2 & R18 KD, 2 images: TGR $-$ AdaBN & +0.013 [-0.002, +0.027] & 0.164 & $\ge 0$ & met \\
R2 & R18 KD, 5 images: TGR $-$ AdaBN & +0.000 [-0.015, +0.015] & 0.958 & none & -- \\
T & R18 KD: TENT $-$ AdaBN & +0.005 [+0.002, +0.009] & 0.007 & within $\pm$0.01, n.s. & partly (within $\pm$0.01, but significant) \\
T & R18 KD: TENT $-$ KD+TGA & -0.012 [-0.025, +0.001] & 0.060 & $<0$, significant & not met (negative, n.s.) \\
T & R18 no-teacher: TENT $-$ AdaBN & +0.005 [+0.002, +0.008] & 0.010 & within $\pm$0.01, n.s. & partly (within $\pm$0.01, but significant) \\
T & R18 no-teacher: TENT $-$ no-teacher adaptation & -0.017 [-0.028, -0.006] & 0.010 & $<0$, significant & met \\
P & R18 KD, $\pi=0.2$: TGR $-$ source positive rate & +0.080 [+0.064, +0.096] & $<$0.001 & $>0$, significant & met \\
P & R18 KD, $\pi=0.6$: TGR $-$ source positive rate & +0.061 [+0.048, +0.077] & $<$0.001 & $>0$, significant & met \\
P & R18 KD, $\pi=0.2$: TGR $-$ source cutoff & +0.038 [+0.020, +0.058] & $<$0.001 & $>0$, significant & met \\
P & R18 KD, $\pi=0.6$: TGR $-$ source cutoff & +0.048 [+0.034, +0.064] & $<$0.001 & $>0$, significant & met \\
P & CNX KD, $\pi=0.2$: TGR $-$ source positive rate & +0.078 [+0.064, +0.094] & $<$0.001 & $>0$, significant & met \\
P & CNX KD, $\pi=0.6$: TGR $-$ source positive rate & +0.083 [+0.067, +0.103] & $<$0.001 & $>0$, significant & met \\
P & CNX KD, $\pi=0.2$: TGR $-$ source cutoff & +0.014 [-0.002, +0.031] & 0.080 & $>0$, significant & not met \\
P & CNX KD, $\pi=0.6$: TGR $-$ source cutoff & +0.051 [+0.038, +0.067] & $<$0.001 & $>0$, significant & met \\
N1a & CNX vs H-optimus-0 teacher: KD+TGA $-$ teacher & -0.009 [-0.022, +0.003] & 0.085 & borderline & NI not established \\
N1b & CNX vs H-optimus-0 teacher: KD+TGR $-$ teacher & -0.013 [-0.026, -0.001] & 0.151 & not non-inferior & NI not established (met) \\
N2c & CNX, Virchow2: KD+TGA $-$ KD (superiority) & +0.044 [+0.031, +0.059] & $<$0.001 & $>0$, significant & met \\
N2d & CNX vs Virchow2 teacher: KD+TGA $-$ teacher & -0.001 [-0.014, +0.012] & 0.004 & as N1a (not NI) & non-inferior (better than expected) \\
N2e & CNX vs Virchow2 teacher: KD+TGR $-$ teacher & -0.016 [-0.031, -0.002] & 0.283 & not non-inferior & NI not established (met) \\
F1 & CNX vs H-optimus-0 LoRA seed 1: KD+TGA $-$ teacher & -0.012 [-0.024, +0.001] & 0.090 & NI not established & NI not established (met) \\
F2 & CNX vs H-optimus-0 LoRA seed 2: KD+TGA $-$ teacher & -0.005 [-0.018, +0.008] & 0.027 & non-inferior & NI not established (not met) \\
F key & CNX vs H-optimus-0, mean over LoRA seeds 0--2: KD+TGA $-$ teacher & -0.009 [-0.020, +0.002] & 0.023 & borderline & non-inferior; partly seen data \\
C1 & CNX, H-optimus-0: KD+TGA $-$ no-teacher adaptation & +0.030 [+0.016, +0.045] & $<$0.001 & 0 to +0.04 & met \\
C2 & CNX, Virchow2: KD+TGA $-$ no-teacher adaptation & +0.030 [+0.017, +0.046] & $<$0.001 & 0 to +0.04 & met \\
\bottomrule
\end{tabular}
}
\end{table}


\FloatBarrier
\section{Implementation details}
\label{app:impl}

\begin{table}[t]
\floatconts{tab:ladder}
{\caption{The repair ladder. Each step acts on one layer of the student and needs different things on the unseen
scanner. All label-free steps are fitted on the unlabelled cells of one half of the new-scanner images and scored on the
other half, then the halves are swapped (cross-fitting).}}
{%
\scriptsize
\setlength{\tabcolsep}{3.5pt}
\begin{tabular}{@{}llcccl@{}}
\toprule
 & & \multicolumn{3}{c}{Needs on the new scanner} & \\
\cmidrule(lr){3-5}
Step & Layer & Unlabelled images & Teacher & Retraining & Students \\
\midrule
No-teacher student (0\%) & -- & -- & -- & -- & both \\
KD, source cutoff & -- & -- & -- & -- & both \\
\quad + TGR & operating point & \checkmark\ (about 5) & probabilities & -- & both \\
\quad + AdaBN & BN statistics & \checkmark & -- & -- & ResNet-18 \\
\quad + TENT & BN scale and shift & \checkmark & -- & at test time & ResNet-18 \\
KD + TGA & representation & \checkmark\ (25 per half) & soft labels & \checkmark & both \\
No-teacher adaptation & representation & \checkmark\ (25 per half) & -- & \checkmark & both \\
Labelled cutoff (upper bound) & operating point & labelled & -- & -- & both \\
Teacher (100\%) & -- & -- & -- & -- & -- \\
\bottomrule
\end{tabular}
}
\end{table}


\paragraph{Data.} Patches are cut around every annotated cell of the MIDOG++ breast subset at the native resolution:
$144\times144$ pixels for the students (random $128\times128$ crop during training, centre crop at test time) and
$256\times256$ pixels for the teachers (centre $224\times224$ crop). Per held-out scanner, the new-scanner set has 1{,}612
(CS2), 1{,}648 (S360) and 1{,}175 (XR) cells. The source validation set consists of eight fixed training images per
source scanner, drawn once with a fixed seed.

\paragraph{Teachers.} Frozen encoders in bfloat16 with LoRA (rank 16, $\alpha=16$, dropout 0.1) on the query--key--value,
output, and both MLP projections of every block; a linear head on [class token, mean patch token]; binary cross-entropy
with positive-class weighting; AdamW (learning rate $10^{-4}$, weight decay $10^{-4}$, one-cycle schedule), batch 16,
5 epochs, dihedral augmentation; the epoch and cutoff with the best source-validation F1 are kept. Normalisation as on
each model card.

\paragraph{Distillation.} For every source cell, the teacher's logit and its feature vector (class token and mean patch
token) are cached for 16 views: the 8 dihedral transforms, each once in the original colours and once with a random
H\&E stain jitter (per-channel scale $1\pm0.15$, shift $\pm0.03$ in the HED space), whose parameters are stored so that
the student sees exactly the same view. Student loss for KD: $T^2\ell(z/T, z_T/T) + 1-\cos(Wf, f_T)$, where $\ell$ is the binary cross-entropy between
the sigmoids of its arguments, $T=2$, $z$ and $z_T$ are the student and teacher logits, $f$ and $f_T$ their features, and
$W$ is a learned linear projection; no expert labels are used. No-teacher student: positive-weighted
BCE with the expert labels on the same views. AdamW (learning rate $3\times10^{-4}$ for ResNet-18, $10^{-4}$ for
ConvNeXt-T; weight decay $10^{-4}$; one-cycle schedule), batch 64, 60 epochs, checkpoint and cutoff by source-validation F1.

\paragraph{TGA and the no-teacher adaptation.} TGA adds the unlabelled cells of the fit half of the new scanner, with
the teacher's cached logits and features on the same 16 views, to the KD training set; nothing else changes. The
no-teacher adaptation keeps the expert labels on source cells and gives the fit-half cells the temperature-softened mean
logit of the ten no-teacher ResNet-18 students after AdaBN on that half, with no feature term. For each student seed and
held-out scanner, two models are trained (fit half A, fit half B), and each scores the other half.

\paragraph{Post-hoc steps.} TGR as in \equationref{eq:tgr}, with the bias found by bounded scalar minimisation in
$[-10, 10]$. AdaBN resets every BatchNorm layer's running statistics and re-estimates them as a cumulative average over
the fit-half cells. TENT trains only the BatchNorm affine parameters with Adam on the mean binary entropy of the fit-half
cells (batch 64, no augmentation), then re-estimates the statistics.

\paragraph{Baselines for the operating point.} Source positive rate: the cutoff on the fit half that reproduces the
student's positive rate on source validation. UPA \citep{roschewitz2023automatic}: piecewise-linear quantile mapping
(101 quantiles) of fit-half scores onto source-validation scores, then the source cutoff. One bias, source mean: the
offset that makes the mean fit-half probability equal the mean source-validation probability. EM
\citep{saerens2002adjusting}, optionally after bias-corrected temperature scaling fitted on labelled source-validation
cells \citep{alexandari2020maximum}, and BBSE \citep{lipton2018detecting} estimate the target prevalence and shift the
logit accordingly. Teacher-referenced variants fit bias and temperature, map quantiles onto the teacher's logits,
reproduce the teacher's positive rate, or choose the F1-optimal cutoff against the teacher's hard labels.

\paragraph{Prevalence shift.} For each held-out scanner and target prevalence $\pi\in\{0.2,\dots,0.6\}$, 20 subsets are
drawn by removing randomly chosen positives (if $\pi$ is below the natural prevalence) or negatives (otherwise); every
method is refitted on each subset with the same cross-fitting, and counts are pooled over the 20 subsets.

\paragraph{Software and hardware.} PyTorch and timm on one NVIDIA RTX 3090 Ti (24 GB). Foundation-model weights were
obtained from their official repositories under their licences.

\FloatBarrier
\section{Additional results}

\FloatBarrier
\subsection{Per-scanner results}
\label{app:perscanner}
\begin{table}[htbp]
\floatconts{tab:app-perscanner}
{\caption{New-scanner F1 per held-out scanner (mean over 10 student seeds) and mean over scanners. The no-teacher rows of the ResNet-18 student do not depend on the teacher and are shown once.}}
{%
\scriptsize
\setlength{\tabcolsep}{3pt}
\begin{tabular}{@{}lcccc@{}}
\toprule
Arm & CS2 & S360 & XR & Mean \\
\midrule
\multicolumn{5}{@{}l}{\textit{ResNet-18, H-optimus-0 teacher}} \\
\quad No-teacher student & 0.754 & 0.635 & 0.745 & 0.711 \\
\quad KD & 0.798 & 0.639 & 0.738 & 0.725 \\
\quad KD + TGR & 0.798 & 0.717 & 0.770 & 0.762 \\
\quad KD + AdaBN + TGR & 0.787 & 0.755 & 0.788 & 0.777 \\
\quad KD + TGA & 0.796 & 0.779 & 0.790 & 0.788 \\
\quad No-teacher adaptation & 0.759 & 0.773 & 0.776 & 0.769 \\
\quad Teacher (seed 0) & 0.833 & 0.806 & 0.803 & 0.814 \\
\multicolumn{5}{@{}l}{\textit{ResNet-18, Virchow2 teacher}} \\
\quad KD & 0.791 & 0.641 & 0.733 & 0.722 \\
\quad KD + TGR & 0.790 & 0.715 & 0.738 & 0.748 \\
\quad KD + AdaBN + TGR & 0.790 & 0.770 & 0.764 & 0.775 \\
\quad KD + TGA & 0.797 & 0.774 & 0.774 & 0.781 \\
\quad Teacher (seed 0) & 0.828 & 0.801 & 0.788 & 0.805 \\
\multicolumn{5}{@{}l}{\textit{ResNet-18, UNI2-h teacher}} \\
\quad KD & 0.776 & 0.606 & 0.731 & 0.704 \\
\quad KD + TGR & 0.781 & 0.681 & 0.769 & 0.744 \\
\quad KD + AdaBN + TGR & 0.779 & 0.742 & 0.784 & 0.768 \\
\quad KD + TGA & 0.775 & 0.756 & 0.783 & 0.771 \\
\quad Teacher (seed 0) & 0.819 & 0.779 & 0.791 & 0.797 \\
\multicolumn{5}{@{}l}{\textit{ResNet-18, UNI teacher}} \\
\quad KD & 0.778 & 0.614 & 0.757 & 0.716 \\
\quad KD + TGR & 0.748 & 0.676 & 0.737 & 0.720 \\
\quad KD + AdaBN + TGR & 0.766 & 0.699 & 0.749 & 0.738 \\
\quad KD + TGA & 0.799 & 0.765 & 0.783 & 0.782 \\
\quad Teacher (seed 0) & 0.786 & 0.756 & 0.758 & 0.767 \\
\multicolumn{5}{@{}l}{\textit{ConvNeXt-T, H-optimus-0 teacher}} \\
\quad No-teacher student & 0.780 & 0.684 & 0.743 & 0.736 \\
\quad KD & 0.813 & 0.714 & 0.768 & 0.765 \\
\quad KD + TGR & 0.828 & 0.785 & 0.788 & 0.801 \\
\quad KD + TGA & 0.803 & 0.803 & 0.809 & 0.805 \\
\quad No-teacher adaptation & 0.759 & 0.777 & 0.789 & 0.775 \\
\quad Teacher (seed 0) & 0.833 & 0.806 & 0.803 & 0.814 \\
\multicolumn{5}{@{}l}{\textit{ConvNeXt-T, Virchow2 teacher}} \\
\quad No-teacher student & 0.780 & 0.684 & 0.743 & 0.736 \\
\quad KD & 0.826 & 0.701 & 0.756 & 0.761 \\
\quad KD + TGR & 0.828 & 0.781 & 0.760 & 0.789 \\
\quad KD + TGA & 0.826 & 0.789 & 0.800 & 0.805 \\
\quad No-teacher adaptation & 0.759 & 0.777 & 0.789 & 0.775 \\
\quad Teacher (seed 0) & 0.828 & 0.801 & 0.788 & 0.805 \\
\bottomrule
\end{tabular}
}
\end{table}


\FloatBarrier
\subsection{The ladder in both orders, the statistics swap, and the UNI teacher}
\label{app:ladder}

\tableref{tab:app-orders} shows why we report Shapley credit: AdaBN alone already moves the operating point, so the
gain of TGR is smaller when it comes second. \tableref{tab:app-bnswap} shows that the retrained (TGA) students' gain on
S360 and XR is not carried by their BatchNorm statistics. The seed-0 UNI teacher chose source cutoffs of 0.25 to 0.28
(seeds 1 and 2: 0.39 to 0.62) and over-calls mitoses on the new scanners; TGR transfers this operating point to the
student (threshold layer $-10\%$ of a small advantage), whereas retraining does not depend on it, so that the retrained
student exceeds this particular teacher ($+0.016$, $+0.004$ to $+0.029$).

\begin{table}[htbp]
\floatconts{tab:app-orders}
{\caption{The two post-hoc layers of the ResNet-18 KD student in both orders (F1 gain of each step) and their order-free Shapley credit with 95\% intervals. AdaBN also moves the operating point, so the threshold gain shrinks when AdaBN comes first.}}
{%
\scriptsize
\setlength{\tabcolsep}{3pt}
\begin{tabular}{@{}lcccccc@{}}
\toprule
Teacher & \makecell{TGR\\first} & \makecell{then\\AdaBN} & \makecell{AdaBN\\first} & \makecell{then\\TGR} & Shapley: threshold & Shapley: statistics \\
\midrule
H-optimus-0 & +0.037 & +0.015 & +0.047 & +0.005 & +0.021 [+0.013, +0.030] & +0.031 [+0.015, +0.049] \\
Virchow2 & +0.026 & +0.027 & +0.047 & +0.005 & +0.016 [+0.008, +0.024] & +0.037 [+0.021, +0.054] \\
UNI2-h & +0.039 & +0.025 & +0.057 & +0.007 & +0.023 [+0.016, +0.032] & +0.041 [+0.023, +0.061] \\
UNI & +0.004 & +0.017 & +0.036 & -0.015 & -0.005 [-0.019, +0.009] & +0.027 [+0.010, +0.044] \\
\bottomrule
\end{tabular}
}
\end{table}

\begin{table}[htbp]
\floatconts{tab:app-bnswap}
{\caption{ResNet-18 students retrained on the new scanner's fit half, evaluated with their own BatchNorm statistics and with statistics re-estimated on source cells only (cutoff re-chosen on source validation). Drop: mean over 10 seeds, interval from a seed bootstrap. On S360 and XR, where retraining helps, the statistics explain little of its gain; on CS2 weights and statistics have co-adapted, so the swap is not interpretable as the statistics' contribution.}}
{%
\scriptsize
\setlength{\tabcolsep}{3pt}
\begin{tabular}{@{}llccc@{}}
\toprule
Student & Scanner & Own statistics & Source statistics & Drop [95\% CI] \\
\midrule
No-teacher adaptation & CS2 & 0.759 & 0.739 & 0.020 [0.014, 0.027] \\
No-teacher adaptation & S360 & 0.773 & 0.769 & 0.004 [-0.001, 0.009] \\
No-teacher adaptation & XR & 0.776 & 0.759 & 0.017 [0.005, 0.027] \\
KD + TGA & CS2 & 0.796 & 0.667 & 0.128 [0.089, 0.170] \\
KD + TGA & S360 & 0.779 & 0.769 & 0.011 [0.002, 0.019] \\
KD + TGA & XR & 0.790 & 0.783 & 0.007 [0.002, 0.013] \\
\bottomrule
\end{tabular}
}
\end{table}


\FloatBarrier
\subsection{Recalibration baselines}
\label{app:baselines}
\begin{table}[t]
\floatconts{tab:recal}
{\caption{Recalibrating the operating point without new-scanner labels (KD students of H-optimus-0). F1 on the unseen
scanner at its natural mitosis prevalence and after resampling its cells to a prevalence of 0.2 or 0.6 (source scanners:
0.32 to 0.45). Bold: best label-free method in the column. $^\ddagger$Also uses the labels of the source-validation
cells. At a single operating point, UPA reduces to keeping the source positive rate. Paired differences to TGR with 95\%
intervals are in \tableref{tab:recal-diff}.}}
{%
\scriptsize
\setlength{\tabcolsep}{4pt}
\begin{tabular}{@{}lccc@{\hspace{10pt}}ccc@{}}
\toprule
 & \multicolumn{3}{c}{ResNet-18 KD student} & \multicolumn{3}{c}{ConvNeXt-T KD student} \\
\cmidrule(lr){2-4}\cmidrule(l){5-7}
Prevalence on the new scanner & natural & 0.2 & 0.6 & natural & 0.2 & 0.6 \\
\midrule
\multicolumn{7}{@{}l}{\textit{No recalibration}} \\
\quad Source cutoff & 0.725 & 0.609 & 0.784 & 0.765 & 0.661 & 0.815 \\
\addlinespace[2pt]
\multicolumn{7}{@{}l}{\textit{Teacher-free (reference: the student's source-validation outputs)}} \\
\quad UPA, quantile alignment \citep{roschewitz2023automatic} & 0.760 & 0.566 & 0.771 & 0.781 & 0.597 & 0.783 \\
\quad Source positive rate & 0.760 & 0.566 & 0.770 & 0.782 & 0.597 & 0.784 \\
\quad One bias, source mean probability & 0.760 & 0.601 & 0.765 & 0.779 & 0.650 & 0.741 \\
\quad EM prior adjustment & 0.552 & 0.411 & 0.705 & 0.671 & 0.489 & 0.793 \\
\quad EM + bias-corrected temperature scaling$^\ddagger$ & 0.521 & 0.407 & 0.681 & 0.678 & 0.497 & 0.794 \\
\quad BBSE$^\ddagger$ & 0.545 & 0.410 & 0.720 & 0.674 & 0.460 & 0.797 \\
\addlinespace[2pt]
\multicolumn{7}{@{}l}{\textit{Teacher-referenced (reference: the teacher on the same unlabelled cells)}} \\
\quad TGR: one bias (ours) & 0.762 & 0.646 & 0.832 & \textbf{0.801} & 0.675 & 0.866 \\
\quad Bias + temperature & 0.760 & \textbf{0.649} & 0.830 & 0.799 & \textbf{0.684} & 0.866 \\
\quad Quantile map onto the teacher & \textbf{0.763} & 0.632 & \textbf{0.839} & 0.799 & 0.669 & 0.867 \\
\quad Teacher positive rate & 0.762 & 0.632 & \textbf{0.839} & 0.799 & 0.668 & 0.867 \\
\quad Cutoff from teacher pseudo-labels & 0.757 & 0.618 & \textbf{0.839} & 0.794 & 0.658 & \textbf{0.870} \\
\addlinespace[2pt]
\multicolumn{7}{@{}l}{\textit{Uses new-scanner labels}} \\
\quad Labelled cutoff (upper bound) & 0.762 & 0.666 & 0.842 & 0.800 & 0.705 & 0.872 \\
\addlinespace[2pt]
\multicolumn{7}{@{}l}{\textit{Reference}} \\
\quad Teacher alone & 0.814 & 0.681 & 0.883 & 0.814 & 0.681 & 0.883 \\
\bottomrule
\end{tabular}
}
\end{table}

\begin{table}[htbp]
\floatconts{tab:recal-diff}
{\caption{TGR minus each recalibration method (F1, 95\% paired-bootstrap interval; descriptive). A separate bootstrap run from \tableref{tab:app-tests}, so shared contrasts can differ in the third decimal of an interval limit.}}
{%
\tiny
\setlength{\tabcolsep}{1.5pt}
\begin{tabular}{@{}lcccccc@{}}
\toprule
Method & \multicolumn{3}{c}{ResNet-18 KD: natural, 0.2, 0.6} & \multicolumn{3}{c}{ConvNeXt-T KD: natural, 0.2, 0.6} \\
\midrule
Source cutoff & \makecell{+0.037\\[-1pt]{\tiny[+0.022, +0.053]}} & \makecell{+0.038\\[-1pt]{\tiny[+0.020, +0.058]}} & \makecell{+0.048\\[-1pt]{\tiny[+0.034, +0.064]}} & \makecell{+0.036\\[-1pt]{\tiny[+0.022, +0.052]}} & \makecell{+0.014\\[-1pt]{\tiny[-0.002, +0.031]}} & \makecell{+0.051\\[-1pt]{\tiny[+0.038, +0.067]}} \\
UPA & \makecell{+0.002\\[-1pt]{\tiny[-0.008, +0.012]}} & \makecell{+0.081\\[-1pt]{\tiny[+0.065, +0.097]}} & \makecell{+0.061\\[-1pt]{\tiny[+0.047, +0.077]}} & \makecell{+0.020\\[-1pt]{\tiny[+0.010, +0.032]}} & \makecell{+0.078\\[-1pt]{\tiny[+0.064, +0.094]}} & \makecell{+0.084\\[-1pt]{\tiny[+0.067, +0.104]}} \\
Source positive rate & \makecell{+0.002\\[-1pt]{\tiny[-0.008, +0.012]}} & \makecell{+0.080\\[-1pt]{\tiny[+0.064, +0.096]}} & \makecell{+0.061\\[-1pt]{\tiny[+0.048, +0.077]}} & \makecell{+0.019\\[-1pt]{\tiny[+0.009, +0.031]}} & \makecell{+0.078\\[-1pt]{\tiny[+0.064, +0.094]}} & \makecell{+0.083\\[-1pt]{\tiny[+0.067, +0.103]}} \\
One bias, source mean & \makecell{+0.001\\[-1pt]{\tiny[-0.007, +0.011]}} & \makecell{+0.046\\[-1pt]{\tiny[+0.031, +0.061]}} & \makecell{+0.067\\[-1pt]{\tiny[+0.050, +0.087]}} & \makecell{+0.022\\[-1pt]{\tiny[+0.010, +0.037]}} & \makecell{+0.025\\[-1pt]{\tiny[+0.016, +0.034]}} & \makecell{+0.126\\[-1pt]{\tiny[+0.096, +0.158]}} \\
EM prior & \makecell{+0.210\\[-1pt]{\tiny[+0.175, +0.240]}} & \makecell{+0.235\\[-1pt]{\tiny[+0.213, +0.257]}} & \makecell{+0.127\\[-1pt]{\tiny[+0.097, +0.158]}} & \makecell{+0.129\\[-1pt]{\tiny[+0.098, +0.161]}} & \makecell{+0.186\\[-1pt]{\tiny[+0.152, +0.216]}} & \makecell{+0.073\\[-1pt]{\tiny[+0.054, +0.095]}} \\
EM + BCTS & \makecell{+0.241\\[-1pt]{\tiny[+0.212, +0.264]}} & \makecell{+0.240\\[-1pt]{\tiny[+0.220, +0.261]}} & \makecell{+0.150\\[-1pt]{\tiny[+0.117, +0.183]}} & \makecell{+0.122\\[-1pt]{\tiny[+0.091, +0.155]}} & \makecell{+0.178\\[-1pt]{\tiny[+0.141, +0.212]}} & \makecell{+0.073\\[-1pt]{\tiny[+0.053, +0.095]}} \\
BBSE & \makecell{+0.216\\[-1pt]{\tiny[+0.180, +0.245]}} & \makecell{+0.236\\[-1pt]{\tiny[+0.213, +0.260]}} & \makecell{+0.112\\[-1pt]{\tiny[+0.084, +0.140]}} & \makecell{+0.127\\[-1pt]{\tiny[+0.096, +0.158]}} & \makecell{+0.215\\[-1pt]{\tiny[+0.177, +0.248]}} & \makecell{+0.070\\[-1pt]{\tiny[+0.050, +0.092]}} \\
Bias + temp. (teacher) & \makecell{+0.002\\[-1pt]{\tiny[+0.000, +0.004]}} & \makecell{-0.002\\[-1pt]{\tiny[-0.007, +0.002]}} & \makecell{+0.001\\[-1pt]{\tiny[+0.001, +0.002]}} & \makecell{+0.001\\[-1pt]{\tiny[-0.001, +0.004]}} & \makecell{-0.009\\[-1pt]{\tiny[-0.013, -0.004]}} & \makecell{+0.001\\[-1pt]{\tiny[-0.000, +0.002]}} \\
Quantile map (teacher) & \makecell{-0.001\\[-1pt]{\tiny[-0.006, +0.004]}} & \makecell{+0.014\\[-1pt]{\tiny[+0.008, +0.021]}} & \makecell{-0.007\\[-1pt]{\tiny[-0.011, -0.004]}} & \makecell{+0.002\\[-1pt]{\tiny[-0.002, +0.006]}} & \makecell{+0.007\\[-1pt]{\tiny[+0.002, +0.013]}} & \makecell{-0.001\\[-1pt]{\tiny[-0.004, +0.003]}} \\
Teacher positive rate & \makecell{-0.001\\[-1pt]{\tiny[-0.006, +0.005]}} & \makecell{+0.014\\[-1pt]{\tiny[+0.008, +0.021]}} & \makecell{-0.007\\[-1pt]{\tiny[-0.010, -0.003]}} & \makecell{+0.002\\[-1pt]{\tiny[-0.002, +0.006]}} & \makecell{+0.007\\[-1pt]{\tiny[+0.002, +0.013]}} & \makecell{-0.000\\[-1pt]{\tiny[-0.003, +0.003]}} \\
Teacher pseudo-label cut. & \makecell{+0.004\\[-1pt]{\tiny[-0.001, +0.011]}} & \makecell{+0.028\\[-1pt]{\tiny[+0.020, +0.037]}} & \makecell{-0.007\\[-1pt]{\tiny[-0.012, -0.003]}} & \makecell{+0.006\\[-1pt]{\tiny[+0.002, +0.012]}} & \makecell{+0.018\\[-1pt]{\tiny[+0.012, +0.025]}} & \makecell{-0.003\\[-1pt]{\tiny[-0.007, +0.001]}} \\
Labelled cutoff & \makecell{-0.001\\[-1pt]{\tiny[-0.007, +0.006]}} & \makecell{-0.020\\[-1pt]{\tiny[-0.032, -0.008]}} & \makecell{-0.010\\[-1pt]{\tiny[-0.016, -0.004]}} & \makecell{+0.000\\[-1pt]{\tiny[-0.006, +0.008]}} & \makecell{-0.030\\[-1pt]{\tiny[-0.046, -0.012]}} & \makecell{-0.006\\[-1pt]{\tiny[-0.011, -0.001]}} \\
Teacher alone & \makecell{-0.052\\[-1pt]{\tiny[-0.069, -0.039]}} & \makecell{-0.035\\[-1pt]{\tiny[-0.056, -0.017]}} & \makecell{-0.052\\[-1pt]{\tiny[-0.066, -0.039]}} & \makecell{-0.013\\[-1pt]{\tiny[-0.026, -0.001]}} & \makecell{-0.006\\[-1pt]{\tiny[-0.025, +0.011]}} & \makecell{-0.017\\[-1pt]{\tiny[-0.027, -0.007]}} \\
\bottomrule
\end{tabular}
}
\end{table}

\begin{table}[htbp]
\floatconts{tab:app-priors}
{\caption{Label-free estimates of the new scanner's mitosis prevalence (mean over 10 student seeds) against the truth. The prior-shift estimators follow the student's own shifted positive rate and amplify it (S360, XR).}}
{%
\scriptsize
\setlength{\tabcolsep}{3pt}
\begin{tabular}{@{}llcccccc@{}}
\toprule
KD student & Scanner & True & \makecell{Student positive\\rate} & EM & EM + BCTS & BBSE & \makecell{Teacher positive\\rate} \\
\midrule
ConvNeXt-T & CS2 & 0.43 & 0.39 & 0.53 & 0.40 & 0.37 & 0.43 \\
ConvNeXt-T & S360 & 0.35 & 0.24 & 0.20 & 0.15 & 0.17 & 0.47 \\
ConvNeXt-T & XR & 0.38 & 0.55 & 0.62 & 0.59 & 0.55 & 0.47 \\
ResNet-18 & CS2 & 0.43 & 0.43 & 0.58 & 0.45 & 0.41 & 0.43 \\
ResNet-18 & S360 & 0.35 & 0.21 & 0.09 & 0.04 & 0.08 & 0.47 \\
ResNet-18 & XR & 0.38 & 0.58 & 0.66 & 0.67 & 0.61 & 0.47 \\
\bottomrule
\end{tabular}
}
\end{table}


\FloatBarrier
\subsection{Few images}
\label{app:few}
\begin{table}[htbp]
\floatconts{tab:app-budget}
{\caption{New-scanner F1 with $n$ images per half from the new scanner (mean over folds, seeds and 10 image draws; the same draws for every method). TGR and AdaBN use the images unlabelled; the labelled cutoff uses their labels. Source cutoff: 0.725 (ResNet-18), 0.765 (ConvNeXt-T).}}
{%
\footnotesize
\setlength{\tabcolsep}{5pt}
\begin{tabular}{@{}cccccc@{}}
\toprule
$n$ & R18: TGR & R18: AdaBN & R18: labelled cutoff & CNX: TGR & CNX: labelled cutoff \\
\midrule
1 & 0.740 & 0.715 & 0.646 & 0.791 & 0.705 \\
2 & 0.755 & 0.743 & 0.734 & 0.798 & 0.775 \\
5 & 0.757 & 0.757 & 0.743 & 0.799 & 0.789 \\
10 & 0.762 & 0.772 & 0.760 & 0.801 & 0.798 \\
25 & 0.762 & 0.772 & 0.762 & 0.801 & 0.800 \\
\bottomrule
\end{tabular}
}
\end{table}


\FloatBarrier
\subsection{Teachers and inference cost}
\label{app:teachers}
\begin{table}[htbp]
\floatconts{tab:app-teachers}
{\caption{Teachers on the unseen scanner: F1 for three LoRA seeds (mean over the three held-out scanners), their mean, and the range of AUROC over seeds; below, seed 0 with eight-view test-time augmentation and a TTA cutoff from source validation. Seed-0 values are re-evaluated from the saved LoRA weights and can differ in the third decimal from the cached predictions used elsewhere (UNI: 0.766 here, 0.767 in \tableref{tab:main}).}}
{%
\footnotesize
\setlength{\tabcolsep}{5pt}
\begin{tabular}{@{}lccccc@{}}
\toprule
Teacher & Seed 0 & Seed 1 & Seed 2 & Mean & AUROC range \\
\midrule
H-optimus-0 & 0.814 & 0.821 & 0.801 & 0.812 & 0.925--0.929 \\
Virchow2 & 0.805 & 0.805 & 0.803 & 0.804 & 0.917--0.918 \\
UNI2-h & 0.797 & 0.792 & 0.787 & 0.792 & 0.913--0.916 \\
UNI & 0.766 & 0.782 & 0.786 & 0.778 & 0.905--0.908 \\
\midrule
H-optimus-0, 8-view TTA (seed 0) & 0.812 & & & & \\
Virchow2, 8-view TTA (seed 0) & 0.797 & & & & \\
\bottomrule
\end{tabular}
}
\end{table}

\begin{table}[htbp]
\floatconts{tab:app-efficiency}
{\caption{Inference cost on one RTX 3090 Ti: random inputs at the resolution used in training, mixed precision (fp16 students, bf16 teachers), 5 warm-up batches, then 30 timed batches, median of 3 repetitions. Parameters include LoRA and heads.}}
{%
\footnotesize
\setlength{\tabcolsep}{5pt}
\begin{tabular}{@{}lrrrrr@{}}
\toprule
Model & Params (M) & Input (px) & Batch & Cells/s & Peak memory (GB) \\
\midrule
ResNet-18 student & 11.2 & 128 & 256 & 23,073 & 0.40 \\
ConvNeXt-T student & 27.8 & 128 & 256 & 7,899 & 0.72 \\
H-optimus-0 teacher & 1,150.5 & 224 & 64 & 80 & 3.29 \\
Virchow2 teacher & 641.7 & 224 & 64 & 128 & 2.15 \\
UNI2-h teacher & 690.8 & 224 & 64 & 132 & 2.41 \\
UNI teacher & 309.6 & 224 & 64 & 342 & 1.08 \\
\bottomrule
\end{tabular}
}
\end{table}


\FloatBarrier
\subsection{TENT and the scanner probe}
\label{app:tent}
\begin{table}[htbp]
\floatconts{tab:app-tent}
{\caption{TENT on the ResNet-18 students (new-scanner F1, mean over seeds and scanners). Primary setting: learning rate $10^{-3}$, one epoch; the others are a descriptive sensitivity analysis.}}
{%
\footnotesize
\setlength{\tabcolsep}{5pt}
\begin{tabular}{@{}lccccc@{}}
\toprule
Student & lr & Epochs & AdaBN & AdaBN + TENT & + TGR \\
\midrule
No-teacher & 0.0001 & 1 & 0.748 & 0.748 & 0.757 \\
No-teacher & 0.001 & 1 & 0.748 & 0.753 & 0.759 \\
No-teacher & 0.001 & 3 & 0.748 & 0.756 & 0.760 \\
No-teacher & 0.01 & 1 & 0.748 & 0.693 & 0.707 \\
KD & 0.0001 & 1 & 0.772 & 0.772 & 0.780 \\
KD & 0.001 & 1 & 0.772 & 0.777 & 0.784 \\
KD & 0.001 & 3 & 0.772 & 0.778 & 0.784 \\
KD & 0.01 & 1 & 0.772 & 0.726 & 0.735 \\
\bottomrule
\end{tabular}
}
\end{table}

\begin{table}[htbp]
\floatconts{tab:app-probe}
{\caption{Linear probe for the training scanner on source-cell features (image-grouped cross-validation, two classes, mean over folds and seeds). All models encode the scanner almost perfectly, so robustness does not mean that scanner information is absent.}}
{%
\footnotesize
\setlength{\tabcolsep}{5pt}
\begin{tabular}{@{}lcc@{}}
\toprule
Features & Dimension & Probe accuracy \\
\midrule
H-optimus-0 & 3072 & 1.000 \\
Virchow2 & 2560 & 0.999 \\
UNI2-h & 3072 & 1.000 \\
UNI & 2048 & 1.000 \\
H-optimus-0 (PCA 512) & 512 & 0.814 \\
Virchow2 (PCA 512) & 512 & 0.824 \\
UNI2-h (PCA 512) & 512 & 0.853 \\
UNI (PCA 512) & 512 & 0.859 \\
ResNet-18 no-teacher & 512 & 0.962 \\
ResNet-18 KD (H-optimus-0) & 512 & 0.999 \\
ResNet-18 KD + TGA (H-optimus-0, half A) & 512 & 0.999 \\
\bottomrule
\end{tabular}
}
\end{table}


\FloatBarrier
\section{Protocols F and C}
\label{app:fc}
Both protocols were frozen together (\tableref{tab:app-protocols}) before any of the models they concern was trained.

\paragraph{Protocol C: no-teacher control for ConvNeXt-T.} The no-teacher adaptation recipe of the ResNet-18 control
was applied unchanged to ConvNeXt-T (60 students: three held-out scanners, two fit halves, ten seeds): expert labels on
the source cells plus the soft answers of the AdaBN-adapted ensemble of ten no-teacher ResNet-18 students on the
unlabelled fit half, scored on the other half. Both pre-specified contrasts (two-sided, Holm) are positive and inside
the expected range of 0 to +0.04: KD + TGA minus no-teacher adaptation is +0.030 (+0.016 to +0.045) with H-optimus-0
and +0.030 (+0.017 to +0.046) with Virchow2 (\tableref{tab:app-tests}). The protocol fixed the wording of the main
text in advance by the share of TGA's gain over the no-teacher student that the no-teacher adaptation reaches: if it
was at least 0.5 for both teachers, the text would say that for both students most of the retraining gain comes from
the unlabelled images. It is 0.57 (0.34 to 0.76) for both teachers, so this wording is used; its interval includes values
below 0.5. The adapted no-teacher ConvNeXt-T reaches 0.759 (CS2), 0.777 (S360) and 0.789 (XR), and stays below the KD
student with TGR ($-0.025$, $-0.040$ to $-0.013$, with H-optimus-0; $-0.014$, $-0.031$ to $-0.000$, with Virchow2;
descriptive).

\paragraph{Protocol F: teacher LoRA seed.} The ConvNeXt-T students of H-optimus-0 were distilled again, with the recipe
unchanged, from LoRA seeds 1 and 2 of the teacher (180 students: KD, and KD + TGA on both fit halves, for ten seeds and
three held-out scanners). The teacher caches were rebuilt exactly as for seed 0, so the stain jitter of the cached views
also differs between teacher seeds, as in a complete re-run of the pipeline. Each student is compared with the
teacher training it was distilled from (one view, the teacher's own source cutoff). The teachers reach 0.814, 0.821 and
0.801 for seeds 0, 1 and 2, the retrained students 0.805, 0.809 and 0.796. The protocol expected the student to land near
0.80 to 0.81 whatever the teacher, so that non-inferiority would fail for seed 1 and hold for seed 2. The first
expectation was met and the second was not: seed 1, $-0.012$ ($-0.024$ to $+0.001$, Holm $p=0.090$); seed 2, $-0.005$
($-0.018$ to $+0.008$), whose interval lies within the margin but whose Holm-adjusted $p$ of 0.027 misses 0.025
(unadjusted 0.013). The pre-specified key secondary estimate, the gap averaged over the three trainings, is $-0.009$
($-0.0198$ to $+0.0023$, one-sided $p=0.023$); it meets the non-inferiority criterion, but only just, and its seed-0 term
was known when the protocol was written. Descriptively, the retrained student recovers 86\% (73\% to 101\%) and 92\%
(74\% to 114\%) of the advantage of the seed-1 and seed-2 teachers (seed 0: 88\%), KD alone 48\% and 63\% (seed 0:
37\%), and KD with TGR stays 0.017 ($-0.029$ to $-0.005$) and 0.007 ($-0.021$ to $+0.006$) below them. Across the three
trainings, the retrained student's F1 varies by 0.013 (0.006 to 0.021), the teacher's by 0.020 (0.010 to 0.030).

\FloatBarrier
\section{Compute, code and data}
\label{app:compute}

The study trained 936 networks (36 LoRA teachers and 900 students) in about 22 GPU hours of training-loop time on
one RTX 3090 Ti (\tableref{tab:app-compute}); 150 further exploratory networks are not reported. MIDOG++ is publicly
available \citep{aubreville2023comprehensive}. Code, per-cell results and the pre-registration protocols are available at \url{https://github.com/AaronG211/scanner-robust-distillation}.

\begin{table}[htbp]
\floatconts{tab:app-compute}
{\caption{Networks trained for the study and their training-loop time on one RTX 3090 Ti (excluding teacher caching and evaluation). Not in the paper: two superseded exploratory variants.}}
{%
\tiny
\setlength{\tabcolsep}{3pt}
\begin{tabular}{@{}llrrc@{}}
\toprule
Group & Networks & $n$ & GPU hours & In paper \\
\midrule
teacher & H-optimus-0 LoRA, seed 0 & 3 & 1.24 & yes \\
teacher & uni-lora, seed 0 & 3 & 0.12 & yes \\
teacher & uni2h-lora, seed 0 & 3 & 0.27 & yes \\
teacher & virchow2-lora, seed 0 & 3 & 0.28 & yes \\
teacher & h-optimus-0-lora, seeds 1-2 & 6 & 1.29 & yes \\
teacher & uni-lora, seeds 1-2 & 6 & 0.24 & yes \\
teacher & uni2h-lora, seeds 1-2 & 6 & 0.54 & yes \\
teacher & virchow2-lora, seeds 1-2 & 6 & 0.56 & yes \\
ResNet-18 & no-teacher student & 30 & 0.26 & yes \\
ResNet-18 & KD, H-optimus-0 & 30 & 0.26 & yes \\
ResNet-18 & KD, Virchow2 / UNI / UNI2-h & 90 & 0.85 & yes \\
ResNet-18 & KD + TGA, H-optimus-0 (halves A/B) & 60 & 0.69 & yes \\
ResNet-18 & KD + TGA, Virchow2 / UNI / UNI2-h (halves A/B) & 180 & 2.31 & yes \\
ResNet-18 & no-teacher adaptation (halves A/B) & 60 & 0.73 & yes \\
ConvNeXt-T & no-teacher student & 30 & 0.66 & yes \\
ConvNeXt-T & KD, H-optimus-0 / Virchow2 & 60 & 1.31 & yes \\
ConvNeXt-T & KD + TGA, H-optimus-0 / Virchow2 (halves A/B) & 120 & 3.55 & yes \\
ConvNeXt-T & no-teacher adaptation (halves A/B; Protocol C) & 60 & 1.73 & yes \\
ConvNeXt-T & KD, H-optimus-0 LoRA seeds 1-2 (Protocol F) & 60 & 1.28 & yes \\
ConvNeXt-T & KD + TGA, H-optimus-0 LoRA seeds 1-2 (halves A/B; Protocol F) & 120 & 3.50 & yes \\
ResNet-18 & KD + target adaptation without cross-fitting (superseded) & 30 & 0.40 & no \\
ResNet-18 & TGA variants (ResNet pseudo-labeller, soft-only) & 120 & 1.48 & no \\
\bottomrule
\end{tabular}
}
\end{table}

\FloatBarrier
